% Source-derived chapter 06: Stratified semi-smallness
% Original source: intersection-complexes-unramified-l-factors
\chapter{Stratified semi-smallness}
\label{intersection-complexes-unramified-l-factors-chapter-06}
\source{intersection-complexes-unramified-l-factors}

\begin{remark}[Source section]
\label{intersection-complexes-unramified-l-factors-chapter-06-source}
\source{intersection-complexes-unramified-l-factors}
\source{intersection-complexes-unramified-l-factors}
This chapter reproduces the corresponding section of the retrieved
LaTeX source, including its definitions, statements, examples, and
proofs. It has not yet been translated into Lean.
\notready
\end{remark}

% The source section title was: Stratified semi-smallness
\label{sect:centralfiber}
\source{intersection-complexes-unramified-l-factors}


We keep the assumptions of \S\ref{sect:Heckeact}, i.e.,  
$B$ acts simply transitively on $X^\circ$ and every simple root of $G$ is
a spherical root of type $T$.  



\smallskip

The main result of this section is the following, which 
will be proved in \S\ref{sect:semismall}. 

\begin{thm}
\source{intersection-complexes-unramified-l-factors}
\notready  \label{thm:barpiperverse}
\source{intersection-complexes-unramified-l-factors}
Under the assumptions above, 
$\bar\pi_!(\IC_{\barY^{\ch\lambda}})$ is perverse
and constructible with respect to the stratification 
$\sA^{\ch\lambda} = \bigcup_{\deg(\mf P)=\ch\lambda} \oo C^{\mf P}$
of Proposition~\ref{prop:Astrata}.
\end{thm}




\subsection{Upper bounds on dimension}

Define 
\begin{align*}
    \barY^{\ch\lambda,\succeq \ch\theta} &= \barY^{\ch\lambda} \xt_{\sM_X} \barM_X^{\ch\theta} \subset \barY^{\ch\lambda} \\
    \olsf Y^{\ch\lambda, \succeq \ch\theta} &= \barY^{\ch\lambda, \succeq\ch\theta} \cap \olsf Y^{\ch\lambda} \subset \olsf S^{\ch\lambda}
\end{align*} 
and let $\sY^{\ch\lambda,\succeq \ch\theta}, \msf Y^{\ch\lambda,\succeq \ch\theta}$ denote the corresponding intersections with $\sY^{\ch\lambda}$
(the notation is  
justified by Proposition~\ref{prop:closure-rel}). 
From the stratification $\barY^{\ch\lambda} = \bigcup {_{\ch\nu}}\barY^{\ch\lambda}$ we deduce that 
\[\olsf Y^{\ch\lambda,\succeq\ch\theta} = \bigcup_{\ch\lambda' \le \ch\lambda} \msf Y^{\ch\lambda',\succeq\ch\theta}.\]


\begin{prop}
\source{intersection-complexes-unramified-l-factors}
\notready \label{prop:centralstrata} 
\source{intersection-complexes-unramified-l-factors}
Let $\ch\lambda \in \mathfrak c_X \sm 0$ and 
$\ch\theta \in \mathfrak c_X^-$. 
For any connected component $\mathscr Y$ of 
$\sY^{\ch\lambda,\ch\theta}$, 
we have
\begin{equation} \label{e:semismall-ineq}
\source{intersection-complexes-unramified-l-factors}
    \dim ( \msf Y^{\ch\lambda} \cap \mathscr Y ) 
    \le \frac 1 2 ( \dim (\mathscr Y) - 1 ). 
\end{equation}

Moreover, whenever the inequality above is an equality for an irreducible component $\msf Y$ of $\msf Y^{\ch\lambda} \cap \mathscr Y$, the same holds for the inequality of Proposition \ref{prop:intersections}; that is, the closure $\olsf Y$ in the affine Grassmannian $\Gr_G$ meets a semi-infinite orbit $\msf S^{\ch\lambda'}$ with $\left< \rho_G, \ch\lambda-\ch\lambda'\right> = \dim \olsf Y$.
\end{prop}



\begin{proof}
Let $\ol{\mathscr Y}$ denote the closure of $\mathscr Y$ in $\barY^{\ch\lambda}$.
Recall from Corollary~\ref{cor:Ycomp} that $\mathscr Y = 
\sY^{\ch\lambda} \times_{\sM_X} \mathscr M$ 
where $\mathscr M$ is a connected component of 
$\sM_X^{\ch\theta}$. 
Then we have $\ol{\mathscr Y} \subset \barY^{\ch\lambda} \xt_{\sM_X} \ol{\mathscr M}$
where $\ol{\mathscr M}$ is the closure of $\mathscr M$ in $\barM_X^{\ch\theta}$.



Fix an irreducible component $\msf Y$ of 
$\msf Y^{\ch\lambda}\cap \mathscr Y$ and let $\olsf Y$ be its closure in 
$\olsf Y^{\ch\lambda} \cap \ol{\mathscr Y}$. 
From Proposition~\ref{prop:intersections} we know, by intersecting $\olsf Y$ with semi-infinite orbits in the affine Grassmannian, that there is a $\ch\lambda'\le \ch\lambda$ with $\msf Y':= \olsf Y \cap \msf S^{\ch\lambda'}$ nonempty of dimension zero, and $d:=\dim \msf Y \le \brac{ \rho_G, \ch\lambda-\ch\lambda'}$. 
Note that $\msf Y' \subset \sY^{\ch\lambda'} \cap \ol{\mathscr Y}
\subset \sY^{\ch\lambda'} \xt_{\sM_X} \ol{\mathscr M}$.


Since $\ch\lambda - \ch\lambda' \in \ch\Lambda^\pos_G$, we can decompose 
$\ch\lambda - \ch\lambda' = \sum_{\alpha\in \Delta_G} n_\alpha (\ch\nu_{D^+_\alpha}+\ch\nu_{D^-_\alpha})$
where the sum is over simple roots and 
$\sum n_\alpha = \brac{\rho_G,\ch\lambda-\ch\lambda'}$. 
By the graded factorization property and Lemma~\ref{lem:Yraise}, there is an \'etale map 
\begin{equation} \label{e:YaddC2}
\source{intersection-complexes-unramified-l-factors}
    (\sY^{\ch\lambda'} \xt_{\sM_X} \ol{\mathscr M}) \oo\xt 
    \underset{\alpha \in \Delta_G}{\oo\prod} (C^{(n_\alpha)}\oo\xt C^{(n_\alpha)}) 
    \to
    \Scr Y^{\ch\lambda} \xt_{\sM_X} \ol{\mathscr M}
\end{equation}
over $\ol{\mathscr M}$. 

If $\ch\lambda' \ne 0$, then $\sY^{\ch\lambda'} \xt_{\sM_X} \ol{\mathscr M}$
contains $\msf Y' \tilde\times C$ and hence 
\eqref{e:YaddC2} implies that 
$\dim (\mathscr Y) \ge 1+2\brac{\rho_G,\ch\lambda-\ch\lambda'}$.
Therefore, 
\begin{equation} \label{e:MVineq}
\source{intersection-complexes-unramified-l-factors}
d\le     \brac{\rho_G,\ch\lambda-\ch\lambda'} \le \frac 1 2 (\dim (\mathscr Y) - 1)
\end{equation}
if $\ch\lambda'\ne 0$. 
This proves the
claim in the case $\ch\lambda' \ne 0$. 

There remains to consider when $d= \brac{\rho_G,\ch\lambda-\ch\lambda'}$ and $\ch\lambda'=0$. 
In particular, $\msf Y^0 \xt_{\sM_X} \barM^{\ch\theta}_X$ is nonempty, which can only
be the case if $\ch\theta = 0$ since $\msf Y^0 = \msf Y^{0,0}=\pt$. 
When considering $\sY^{\ch\lambda,0}$, we may
use \S\ref{sect:freemonoid} to reduce to assuming that $\mf c_X = \mbb N^{\Cal D}$. 
In this case again by Proposition \ref{prop:intersections}, there is a simple root $\alpha$ 
and an irreducible component $\msf Y_{d-1}$ of $\olsf Y \cap \msf S^{\ch\alpha}$
of dimension $1$ such that $\ol{\msf Y_{d-1}}$ contains $\msf Y^{0,0} = \pt$. 
This implies that $\msf Y_{d-1}$ is an irreducible component of 
$\msf Y^{\ch\alpha,0}$ of dimension $1$. 
Since $\mf c_X$ is now the free monoid, $\ch\alpha = \ch\nu_{D_\alpha^+} + \ch\nu_{D_\alpha^-}$
for $\Cal D(\alpha) = \{ D_\alpha^+, D_\alpha^- \}$. 
Lemma~\ref{lem:Ycolor} implies that $\msf Y^{\ch\alpha,0}$ is empty, a contradiction.
\end{proof}



\subsection{Connected components of open Zastava} \label{sect:connected-open}
\source{intersection-complexes-unramified-l-factors}
Recall from \S\ref{sect:Y0cc} that $\sY_X^{?,0}=\sY_{X^\bullet}$ is a disjoint union of subschemes
$\sY_{X^\bullet}^{D}$, indexed by $D \in \mbb N^{\Cal D}$, i.e., multisets of colors. 
Therefore, we have a map $\pi_0(\sY_{X^\bullet}) \to \mbb N^{\Cal D}$. 
On the other hand, Corollary~\ref{cor:compY} gives an injection 
$\pi_0(\sY_{X^\bullet}) \into \pi_1(H) \times \mf c_X$ 
where $^{\ch\mu}\sY^{\ch\lambda,0} \mapsto (\ch\mu,\ch\lambda)$ when the former
is nonempty. 

Observe that $\pi_1(H) \ot \bbQ \cong \mathcal X(H) \ot \bbQ$. 
Then tensoring \eqref{e:color-ses} by $\bbQ$, we get an injection 
\begin{equation} \label{e:color-inj-piH}
\source{intersection-complexes-unramified-l-factors}
    \mathbb Z^{\Cal D} \into \mbb Q^{\Cal D} \into (\pi_1(H) \xt \ch\Lambda_X) \ot \bbQ. 
\end{equation}
One can check that the maps above fit into a commutative diagram 
\begin{equation} \label{e:Y0toBunH}
\source{intersection-complexes-unramified-l-factors}
\begin{tikzcd}
\pi_0(\sY_{X^\bullet}) \ar[r, hook] \ar[d] & \pi_1(H)\times \mf c_X \ar[d] \\ 
\mbb N^{\Cal D} \ar[r, hook, "{\eqref{e:color-inj-piH}}"] & 
(\pi_1(H)\times \ch\Lambda_X) \ot \bbQ 
\end{tikzcd}
\end{equation}
We now show that the left vertical arrow is a bijection. 





\begin{lem}
\source{intersection-complexes-unramified-l-factors}
\notready \label{lem:oYconnected}
\source{intersection-complexes-unramified-l-factors}
For $D \in \mathbb N^{\Cal D}$, 
the smooth scheme $\sY^{D}_{X^\bullet}$ is connected of dimension $\len(D)$.
\end{lem}
\begin{proof}
As discussed in \S\ref{sect:Y0cc}, we may assume that $X=X^\can$ and 
$\mathfrak c_X = \mathbb N^{\Cal D}$. 
Now the claim is equivalent to showing that, under those assumptions, 
$\sY^{\ch\lambda,0}$ is connected for $\ch\lambda = \sum n_{D'} \ch\nu_{D'}$. 
One deduces from the graded factorization property of $\sY$ that 
if $\sY^{\ch\lambda,0}$ is not connected, there must exist a possibly different $\ch\lambda$
with a connected component $\mathscr Y$ contained entirely in the preimage of the diagonal 
$\sY^{\ch\lambda,0} \xt_{\sA^{\ch\lambda}, \delta^{\ch\lambda}} C$. 
Then $\dim (\msf Y^{\ch\lambda} \cap \mathscr Y) = \dim \mathscr Y - 1$,
and the dimension inequality of Proposition~\ref{prop:centralstrata} can only hold if
$\dim \mathscr Y = 1$.
By Corollary~\ref{cor:Ycomp}, the component $\mathscr Y$ must be of the form 
\[ \mathscr Y = {^{\ch\mu}\sY^{\ch\lambda,0}} = 
\Bun_H^{\ch\mu} \xt_{\Bun_G} \Bun_B^{-\ch\lambda} \]
for some $\ch\mu\in \pi_1(H)$. 
Now by the graded factorization property and Lemma~\ref{lem:Yraise}, 
we have an \'etale map 
\[ ^{\ch\mu}\sY^{\ch\lambda,0} \oo\xt 
\prod^\circ_{\alpha\in\Delta_G} C^{(N_{\alpha})} \oo\xt C^{(N_{\alpha})} \to 
{^{\ch\mu}}\sY^{\ch\lambda',0} \]
for $\ch\lambda' = \ch\lambda + \sum N_\alpha\ch\alpha$ and any $N_\alpha$ large enough. 
Thus, $\dim {^{\ch\mu}}\sY^{\ch\lambda',0} = 1+2\sum N_\alpha$. 
By Lemma~\ref{lem:BunBtoGsm} we may assume that 
$\Bun_B^{-\ch\lambda'} \to \Bun_G$ is smooth. 
Note that $\dim \Bun_H^{\ch\mu}$ only depends on the image of $\ch\mu$ in $\pi_1(H)\ot \bbQ$. 
On the other hand, the commutative diagram \eqref{e:Y0toBunH} 
says that this image is determined by $\ch\lambda'$. 
These observations imply that $\ds \sY^{\ch\lambda',0}
= \Bun_H \xt_{\Bun_G} \Bun_B^{-\ch\lambda'}$ is equidimensional. 
However we know that $\sY^{\ch\lambda',0}$ has a connected component 
birational to $\sA^{\ch\lambda'}$, which is of dimension $\sum n_{D'} + 2\sum N_\alpha$.
The equality 
\[ \textstyle \dim {^{\ch\mu}}\sY^{\ch\lambda',0} = 1+2\sum N_\alpha = \sum n_{D'} + 2\sum N_\alpha \]
forces $\ch\lambda = \ch\nu_{D'}$ for some color $D'\in \Cal D$, hence $D=D'$.
Lemma~\ref{lem:Ycolor} now implies that $\sY^{\ch\nu_D,0}$ is connected.
\end{proof}

\begin{cor}
\source{intersection-complexes-unramified-l-factors}
\notready \label{cor:Ycomp2}
\source{intersection-complexes-unramified-l-factors}
For every $\ch\lambda\in\mf c_X$, $\ch\theta\in\mf c_X^-$, the connected components of $\sY^{\ch\lambda,\ch\theta}$ are
in bijection with the closures of 
\[ \sY^{D}_{X^\bullet} \oo\xt \sY^{\ch\theta,\ch\theta} \into \sY^{\ch\lambda,\ch\theta} \]
for $D\in \mbb N^{\Cal D}$ such that $\varrho_X(D) = \ch\lambda-\ch\theta$. 
\end{cor}
\begin{proof}
Immediate from Lemmas~\ref{lem:comp-cover}(i) and \ref{lem:oYconnected}.
\end{proof}


\subsection{Stratified semi-smallness} \label{sect:semismall}
\source{intersection-complexes-unramified-l-factors}

Following \cite[\S 4]{MV}, we will use  
the notion of a stratified semi-small map, which we now review.
Let $f : Y \to A$ be a proper map between two stratified spaces 
$(Y,\Scr S)$ and $(A,\Scr T)$. Suppose all strata 
are smooth and connected and each $f(S), S\in \Scr S$ is a union of strata 
$S' \in \Scr T$. We say $f$ is \'etale-\emph{locally trivial} 
(in the stratified sense) if whenever $S' \subset f(S)$, the restriction of 
$f$ to $S\cap f^{-1}(S') \to S'$ is \'etale-locally
a trivial fibration. We say that $f$ is \emph{stratified semi-small} 
if it is \'etale-locally trivial and 
for any $S\in \Scr S$ and any $S'\in \Scr T$ such that
$S' \subset f(S)$ we have 
\begin{equation} \label{e:ss-ineq}
\source{intersection-complexes-unramified-l-factors}
    \dim ( f^{-1}(a) \cap S) \le \frac 1 2 (\dim S - \dim S') 
\end{equation}
for any (and thus all) $a\in S'$. 

The notion of stratified semi-smallness is relevant due to 
the observation below, which follows from dimension counting and the 
definition of the perverse t-structure:

\begin{lem}[{\cite[Lemma 4.3]{MV}}]
\source{intersection-complexes-unramified-l-factors}
\notready  \label{lem:semismallperverse}
\source{intersection-complexes-unramified-l-factors}
If $f$ is a stratified semi-small map then $f_*(\Scr F) \in \rmP_{\Scr T}(A)$ for all 
$\Scr F \in \rmP_{\Scr S}(Y)$. 
\end{lem}

Note that the lemma holds even if the stratifications are not Whitney. 
In this case we simply define $\rmP_{\Scr S}(Y) := \rmP(Y) \cap \rmD^b_{\Scr S}(Y)$
to be the subcategory of perverse sheaves that are $\Scr S$-constructible, i.e., 
the $\Scr F \in \rmP(Y)$ such that $\rmH^i(\Scr F)|_S$ is a local system
of finite rank for all $i\in \bbZ$ and $S\in \Scr S$. 

\subsubsection{}
Let us return to our situation: consider the proper map 
$\bar\pi : \barY^{\ch\lambda} \to \sA^{\ch\lambda}$.

We have the smooth stratification defined in Proposition~\ref{prop:barYstrata}, 
\[ \barY^{\ch\lambda} = \bigcup_{\ch\nu,\ch\Theta} 
{}_{\ch\nu}\barY^{\ch\lambda,\ch\Theta}, \qquad 
{_{\ch\nu}}\barY^{\ch\lambda,\ch\Theta} \cong 
C_{\ch\nu} \xt \sY^{\ch\lambda-\ch\nu,\ch\Theta} \] 
for $\ch\nu \in \ch\Lambda^\pos_G,\, 
\ch\Theta \in \Sym^\infty(\mathfrak c_X^-\sm 0)$.
Let $(\barY^{\ch\lambda},\Scr S)$ denote the stratification by
 the connected components of
$_{\ch\nu}\sY^{\ch\lambda, \ch\Theta}$.
We will not show that $\Scr S$ is a Whitney stratification; for our purposes 
the following suffices: 

\begin{lem}
\source{intersection-complexes-unramified-l-factors}
\notready \label{lem:barICwhitney}
\source{intersection-complexes-unramified-l-factors}
For any $\ch\lambda \in \mathfrak c_X$, 
the IC complex of $\barY^{\ch\lambda}$ is $\Scr S$-constructible, 
i.e., 
\[ \IC_{\barY^{\ch\lambda}} \in \rmP_{\Scr S}(\barY^{\ch\lambda}). \]
\end{lem}
\begin{proof} 
Let $\ch\nu \in \ch\Lambda^\pos_G$ be such that 
$\barY^{\ch\lambda} = {_{\le\ch\nu}}\barY^{\ch\lambda}$, which exists 
since $\barY^{\ch\lambda}$ is of finite type. 
For any $\ch\mu \in \ch\Lambda^\pos_G$ large enough, we have
a smooth correspondence preserving stratifications
\[ \barY^{\ch\lambda} \leftarrow \barY^{\ch\lambda}\oo\xt \sY^{\ch\mu,0} \to 
{_{\le\ch\nu}}\barY^{\ch\lambda+\ch\mu} =: Y\]
where the right arrow comes from the graded factorization property of $\barY$.
Thus, it suffices to check that 
$\IC_Y$ is $\Scr S$-constructible. 
By \eqref{e:ICbar=bt}, we have an isomorphism 
\[ \IC_Y \cong \Bigl.\Bigl(\IC_{\sM_X} \bt_{\Bun_G} \IC_{\ol\Bun_B^{-\ch\lambda-\ch\mu}}\Bigr)\Bigr|_Y. \] 
Now Proposition~\ref{prop:globwhitney} implies that 
$\IC_{\sM_X}$ is constructible with respect to the fine stratification on 
$\sM_X$, and Theorem~\ref{thm:BFGM} implies that $\IC_{\ol\Bun_B}$ is constructible
with respect to the stratification by defect. 
Thus, it follows that $\IC_Y$ is $\Scr S$-constructible.
\end{proof}


Let $(\Scr A^{\ch\lambda}, \Scr T)$ denote the stratification 
from Proposition~\ref{prop:Astrata}, 
\[ \Scr A^{\ch\lambda} = \bigcup_{\mf P} \oo C^{\mf P} \]
for $\mf P\in \Sym^\infty(\mathfrak c_X\sm 0)$ such that 
$\deg(\mf P)=\ch\lambda$.  





Now Theorem~\ref{thm:barpiperverse} follows 
from Lemmas~\ref{lem:semismallperverse}, \ref{lem:barICwhitney} 
and the following theorem:

\begin{thm}
\source{intersection-complexes-unramified-l-factors}
\notready \label{thm:semismall}
\source{intersection-complexes-unramified-l-factors}
The map $\bar\pi : (\barY^{\ch\lambda}, \Scr S) \to (\Scr A^{\ch\lambda}, \Scr T)$ is stratified semi-small. 
\end{thm}
\begin{proof}
Fix $\mf P \in \Sym^\infty(\mathfrak c_X\sm 0)$ such that 
$\ch\lambda=\deg(\mf P)$. 
Let 
$I = \{ 1,\dotsc, \abs{\mf P}\}$ be the finite set of cardinality 
$\abs{\mf P}$. We fix an ordering 
$\mf P = \sum_{i\in I} [\ch\lambda_i]$ on our partition $\mf P$,
so each $\ch\lambda_i \in \mathfrak c_X$. 
Let $\oo C^I$ denote the $I$-fold product $C^I$ with the diagonal
divisor removed. Then the map $\oo C^I \to \oo C^{\mf P}$
corresponds to 
choosing an ordering on an unordered multiset of points in $C$.

We leave it to the reader to check that $\bar\pi(S)$ is a union of
strata for $S \in \Scr S$. 
It is also a consequence of the graded factorization property
that for 
$\ch\nu\in \ch\Lambda^\pos_G, \ch\Theta \in \Sym^\infty(\mathfrak c_X \sm 0)$,
the fiber product 
\begin{equation} \label{e:preimagefactorize}
\source{intersection-complexes-unramified-l-factors}
    _{\ch\nu}\barY^{\ch\lambda,\ch\Theta} 
    \xt_{\sA^{\ch\lambda}} \oo C^I  
    =   \bigsqcup_{i \mapsto \ch\nu_i,\ch\lambda'_i, \ch\theta_i}
    \underset{i\in I}{\oo\prod} (\msf Y^{\ch\lambda'_i, \ch\theta_i} \ttimes C ) 
\end{equation}
is equal to a disjoint union of open and closed subschemes,
running over all assignments $i\in I \mapsto \ch\nu_i \in \ch\Lambda^\pos_G, 
\ch\lambda'_i \in \mf c_X, \ch\theta_i \in \mf c_X^-$ (including zero)
such that $\ch\nu_i + \ch\lambda'_i = \ch\lambda_i$
and $\sum_I [\ch\theta_i] = \ch\Theta + (\abs I - \abs{\ch\Theta})[0] \in 
\Sym^\infty(\mf c_X)$ as a multiset \emph{with zero}. 
 

Since $\oo C^I \to \oo C^{\mf P}$ is finite \'etale,
(and everything is of finite type) we deduce that the restriction of $\bar\pi$ to 
$_{\ch\nu}\barY^{\ch\lambda,\ch\Theta} \cap \bar\pi^{-1}( \oo C^{\mf P} )
    \to \oo C^{\mf P}$
is \'etale-locally a trivial fibration.

It remains to check the dimension inequality \eqref{e:ss-ineq}.
Let $\mf P, \ch\nu, \ch\Theta$ be as before, and  
fix a connected component $S$ of $_{\ch\nu}\barY^{\ch\lambda,\ch\Theta}$.
We deduce from 
\eqref{e:preimagefactorize} that for any $a \in \oo C^{\mf P}$,
the restricted fiber 
$\bar\pi^{-1}(a) \cap S$
is contained in a union of 
$\oo\prod_I (\msf Y^{\ch\lambda'_i,\ch\theta_i}\cap S_i)$ for 
$\ch\nu_i,\ch\lambda'_i,\ch\theta_i$ as above
and $S_i$ some connected component of $\sY^{\ch\lambda'_i,\ch\theta_i}$. 
By Proposition~\ref{prop:centralstrata}, we have 
\begin{equation}  \label{e:dimYile}
\source{intersection-complexes-unramified-l-factors}
\dim( \msf Y^{\ch\lambda'_i, \ch\theta_i} \cap S_i ) \le 
\frac 1 2 ( \dim(S_i) - 1 ).
\end{equation}
The image of the composition
\[ \underset{i\in I}{\oo\prod} (C_{\nu_i} \xt S_i) \to \underset{i\in I}{\oo\prod}
{_{\ch\nu_i} \barY^{\ch\lambda_i,\ch\theta_i} } \to {_{\ch\nu}\barY^{\ch\lambda,\ch\Theta}} \]
is connected, so it must be contained in $S$. 
Thus, $\dim(S) \ge \brac{\rho_G,\ch\nu} + \sum_I \dim(S_i)$.
Summing \eqref{e:dimYile} over $I$, we get 
\[ \dim( \bar\pi^{-1}(a) \cap S ) \le 
\frac 1 2 ( \dim(S) - \dim( \oo C^{\mf P} )- \brac{\rho_G,\ch\nu}  ), \]
which establishes the inequality \eqref{e:ss-ineq}.
\end{proof}

\subsection{Euler product}
Let us explain how we can combine the graded 
factorization property of $\barY$ with 
Theorem~\ref{thm:semismall} to deduce that 
$\bar\pi_!(\IC_{\barY})$ ``looks like an Euler product''. 

In the expression that we are about to obtain, a special role will be played by those strata $\sY^{\ch\lambda,\ch\theta}$ of $\sY^{\ch\lambda}$ which correspond to elements $\ch\theta\in \Cal D^G_{\mathrm{sat}}(X) \cup \{0\}$. Recall, by Corollary \ref{cor:Ycomp}, that the closures of those strata are unions of irreducible components of $\sY$. Let 
\begin{equation}\label{e:crystallambda} \mB_{X,\ch\lambda} = \bigcup_{\ch\theta\in \Cal D^G_{\mathrm{sat}}(X) \cup \{0\}}\bigcup_{\mathscr Y} {\mB_{\mathscr Y}},
\source{intersection-complexes-unramified-l-factors}
\end{equation}
where $\mathscr Y$ runs over all irreducible components of $\sY^{\ch\lambda,\ch\theta}$, and ${\mB_{\mathscr Y}}$ is the set of those irreducible components $\msf b$ of the central fiber $\msf Y^{\ch\lambda}\cap \mathscr Y$ for which the inequality of \eqref{e:semismall-ineq} is an equality, that is, 
\begin{equation}
 \dim\msf b = \frac{1}{2} (\dim \mathscr Y-1).
\end{equation}
Such components will be said to be of \emph{critical dimension}; we will explicate this dimension in Proposition \ref{prop:Vbasis}. The sets $\mB_{X,\ch\lambda}$ will define the \emph{crystal of $X$} in Section \ref{sect:crystal}. For now, we treat them as a black box.


We let $V_{X,\ch\lambda}$ denote the free vector space on $\mB_{X,\ch\lambda}$, that is, 
\begin{equation}\label{e:defVX}
\source{intersection-complexes-unramified-l-factors}
V_{X,\ch\lambda} = \bigoplus_{\mB_{X,\ch\lambda}} \ol\bbQ_\ell.
\end{equation}


Note that, when $X$ is defined over a finite field $\mbb F$, and $k$ is its algebraic closure, the (geometric) Frobenius morphism induces a dimension-preserving bijection between the sets $\mB_{X,\ch\lambda}$ and $\mB_{X,\Fr\ch\lambda}$, for every $\ch\lambda\in\ch\Lambda$. Hence, $\Fr$ acts naturally on the sum of vector spaces $\bigoplus_{\ch\lambda\in\mathfrak c_X} V_{X,\ch\lambda}$. 

For a partition $\mf R = \sum_{\ch\mu \in \mathfrak c_X \sm 0} N_{\ch\mu} [\ch\mu] \in \Sym^\infty(\mathfrak c_X\sm 0)$, 
let $\iota^{\mf R} : \oo C^{\mf R} := \oo\prod \oo C^{(N_{\ch\mu})}
\into \sA^{\ch\lambda}$ denote the locally closed embedding, with $\ch\lambda = \deg(\mf R)$.
This extends to a finite map $\bar\iota^{\mf R} : C^{\mf R}:=\prod C^{(N_{\ch\mu})} \to \sA^{\ch\lambda}$ which is the normalization of the closure of $\oo C^{\mf R}$ in 
$\sA^{\ch\lambda}$.

\begin{prop}
\source{intersection-complexes-unramified-l-factors}
\notready \label{prop:format}
\source{intersection-complexes-unramified-l-factors}
For $\ch\lambda \in \mathfrak c_X$, there
exists a canonical isomorphism 
\begin{equation} \label{e:format1}
\source{intersection-complexes-unramified-l-factors}
 \bar\pi_!(\IC_{\barY^{\ch\lambda}}) \cong 
\bigoplus_{\deg(\mf R)=\ch\lambda} 
\Bigl( \bigotimes_{\ch\mu} 
\Sym^{N_{\ch\mu}}(V_{X,\ch\mu}) \Bigr) \ot 
\bar\iota^{\mf R}_!( \IC_{C^{\mf R}} )
\end{equation}
where $\mf R = \sum_{\ch\mu \in \mathfrak c_X\sm 0} 
N_{\ch\mu} [\ch\mu]$ and the spaces $V_{X,\ch\mu}$ are defined by \eqref{e:defVX}. 

When $X$ is defined over a finite field $\mbb F$, and $k$ is its algebraic closure, this isomorphism is Galois-equivariant. 
\end{prop}

The implied Galois action on the right hand side  of \eqref{e:format1} is the one obtained by the action of Frobenius on the sum of spaces $V_{X,\ch\mu}$ (by permuting their basis elements), and the standard Weil structure $\ol\bbQ_\ell(\frac{|\mf R|}{2})[|\mf R|]$ on $\IC_{C^{\mf R}}$.

Note that if $\mf R = [\ch\lambda]$ is the singleton partition, then
$C^{[\ch\lambda]}=C$ and 
$\iota^{[\ch\lambda]}=\delta^{\ch\lambda} : C \into \sA^{\ch\lambda}$
is the diagonal embedding. 
The corresponding summand of $\bar\pi_!(\IC_{\barY^{\ch\lambda}})$
above is $V_{X,\ch\lambda} \ot \IC_C$. 
We call this the \emph{diagonal contribution} of  
$\bar\pi_!(\IC_{\barY^{\ch\lambda}})$. 

\begin{proof}
The proof follows the same logic as \cite[\S 5.4, 5.11]{BFGM}. 
Theorem~\ref{thm:barpiperverse} implies that 
$\bar\pi_!(\IC_{\barY^{\ch\lambda}})$ is perverse, and 
the decomposition theorem (\cite[Th\'eor\`eme 6.2.5]{BBDG}) 
implies that it is semisimple. 
Since $\bar\pi_!(\IC_{\barY^{\ch\lambda}})$ is constructible
with respect to the stratification by $\iota^{\mf R} : \oo C^{\mf R} \into \sA^{\ch\lambda}$ for $\mf R \in \Sym^\infty(\mathfrak c_X\sm 0),\,
\deg(\mf R)=\ch\lambda$, we deduce that there exists a \emph{canonical} 
decomposition
\begin{equation}\label{e:format2} 
\source{intersection-complexes-unramified-l-factors}
\bar\pi_!(\IC_{\barY^{\ch\lambda}}) \cong 
\bigoplus_{\deg(\mf R)=\ch\lambda} \iota^{\mf R}_{!*}(\Scr L^{\mf R}) [\abs{\mf R}] 
\end{equation}
where $\iota^{\mf R}_{!*}$ denotes the middle extension functor
along the locally closed embedding, and $\Scr L^{\mf R}$ 
is a local system on $\oo C^{\mf R}$.

Now consider the singleton partition $[\ch\lambda]$. 
For $v \in \abs C$ let $\delta^{\ch\lambda}_v : v \to\sA^{\ch\lambda}$
denote the composition of $v\to C$ with $\delta^{\ch\lambda}:C\to \sA^{\ch\lambda}$.
Recall from \S\ref{sect:YttimesC} 
that $\barY^{\ch\lambda} \xt_{\sA^{\ch\lambda}, \delta^{\ch\lambda}} C
\cong \olsf Y^{\ch\lambda} \xt^{\Aut k\tbrac t} C^\wedge$. 
Taking the $*$-pullback along $\delta_v^{\ch\lambda}$ of \eqref{e:format2}, we 
have 
\[ R\Gamma( \olsf Y^{\ch\lambda}, \IC_{\barY^{\ch\lambda}}|^*_{\olsf Y^{\ch\lambda}})
\cong \bigoplus_{\deg(\mf R)=\ch\lambda} (\delta^{\ch\lambda}_v)^* \iota^{\mf R}_{!*}(\Scr L^{\mf R})[\abs{\mf R}]. \]
For $\mf R = [\ch\lambda]$, we have $\Scr L^{[\ch\lambda]}$ is a local system 
on $C$, so $(\delta_v^{\ch\lambda})^* \iota^{[\ch\lambda]}_{!*}(\Scr L^{[\ch\lambda]})[1]$
lives in cohomological degree $-1$. 
For $\mf R \ne [\ch\lambda]$, we have $\dim C^{\mf R} > 1$
so the uniqueness property of middle extension implies that
$(\delta^{\ch\lambda})^* \iota^{\mf R}_{!*}(\Scr L^{\mf R})[\abs{\mf R}]$ 
has lisse cohomology sheaves on $C$ and lives in usual cohomological degrees $< -1$. 
Therefore, we deduce that 
\[ 
    \Scr L^{[\ch\lambda]} |^*_{v\to C} = H_c^{-1}(\olsf Y^{\ch\lambda}, \IC_{\barY^{\ch\lambda}}|^*_{\olsf Y^{\ch\lambda}}),
\]
which is the top cohomological degree. 

Recall from Corollary \ref{cor:compY} that the irreducible components of $\sY^{\ch\lambda}$ are naturally parametrized by a subset of $\pi_1(H) \xt \Sym^\infty(\Cal D^G_{\mathrm{sat}}(X))$. 
Let $\sY^{\ch\lambda,\circ}$ denote the union of $\sY^{\ch\lambda,\ch\Theta}$ 
for all $\ch\Theta \in \Sym^\infty(\Cal D^G_{\mathrm{sat}}(X))$. 
Then $\sY^{\ch\lambda,\circ}$ is a disjoint union of smooth connected components
in bijection with the irreducible components of $\sY^{\ch\lambda}$, by Corollary \ref{cor:Ycomp}.   
Again, the uniqueness property of the IC complex implies that
$\IC_{\barY^{\ch\lambda}}|^*_{\barY^{\ch\lambda}\sm \sY^{\ch\lambda,\circ}}$ 
lives in strictly negative perverse cohomological degrees. 
Since $\bar\pi_!$ is perverse $t$-exact by Theorem~\ref{thm:semismall}, we
deduce that 
$\bar\pi_!( \IC_{\barY^{\ch\lambda}}|^*_{\barY^{\ch\lambda}\sm \sY^{\ch\lambda,\circ}} )$ 
is constructible and lives in strictly negative perverse degrees. 
This in turn implies that $(\delta^{\ch\lambda}_v)^* \bar\pi_!( \IC_{\barY^{\ch\lambda}}|^*_{\barY^{\ch\lambda}\sm \sY^{\ch\lambda,\circ}} )$ lives in (usual=perverse) degrees $<-1$. 
We conclude that 
\[
    \Scr L^{[\ch\lambda]}|^*_{v\to C} = 
    H^{-1}_c(\msf Y^{\ch\lambda}, (\IC_{\barY^{\ch\lambda}}|_{\sY^{\ch\lambda,\circ}})|^*_{\msf Y^{\ch\lambda}} ) = 
\bigoplus_{\mathscr Y} H_c^{\dim \mathscr Y - 1}(\msf Y^{\ch\lambda} \cap \mathscr Y, \ol\bbQ_\ell({\textstyle \frac{\dim \mathscr Y}{2}})),\]
with the sum running over all irreducible components of $\sY^{\ch\lambda,\circ}$.
Note that $\msf Y^{\ch\lambda}\cap \sY^{\ch\lambda,\ch\Theta}$ is empty unless 
$\ch\Theta = [\ch\theta]$ is singleton, for $\ch\theta \in \Cal D^G_{\mathrm{sat}}(X) \cup \{0\}$. 
Moreover, the right hand side consists of only \emph{top} cohomological degrees, by Proposition~\ref{prop:centralstrata}, so 
$H_c^{\dim \mathscr Y - 1}(\msf Y^{\ch\lambda} \cap \mathscr Y, \ol\bbQ_\ell (\frac{\dim \mathscr Y}{2}))$ 
is equal to the sum 
\[\bigoplus_{\mf B_{\mathscr Y}} \ol\bbQ_\ell(\smallfrac{1}{2}),\]
where $\mf B_{\mathscr Y}$ is the set of irreducible components of $\msf Y^{\ch\lambda}\cap \mathscr Y$ of dimension $\frac{\dim \mathscr Y - 1}{2}$.  
In particular, $\Scr L^{[\ch\lambda]}|^*_{v\to C}$ has trivial monodromy under $\Aut k\tbrac t$, 
so we  deduce that $\Scr L^{[\ch\lambda]} \cong V_{X,\ch\lambda} \otimes 
\IC_C$ where $V_{X,\ch\lambda}$ is defined by \eqref{e:defVX} (with the $(\frac{1}{2})$-twist absorbed by $\IC_C$). 

Next, consider an arbitrary partition 
$\mf R = \sum_{\ch\mu\in \mathfrak c_X\sm 0} N_{\ch\mu} [\ch\mu]$. 
We defined $\oo C^{\mf R} = \oo\prod_{\ch\mu} \oo C^{(N_{\ch\mu})}$.
By the graded factorization property, we have a diagram
with Cartesian squares
\[ 
\begin{tikzcd}
    \oo\prod_{\ch\mu} (\olsf Y^{\ch\mu}\ttimes C)^{\oo\xt N_{\ch\mu}} \ar[r, hook] \ar[d] &
    \oo\prod_{\ch\mu} (\barY^{\ch\mu})^{\oo\xt N_{\ch\mu}}
    \ar[r, "\text{\'etale}"] \ar[d] & \barY^{\ch\lambda}  
    \ar[d] \\ 
    \oo\prod \oo C^{N_{\ch\mu}} \ar[r, hook] &
    \oo\prod_{\ch\mu} (\sA^{\ch\mu})^{\oo\xt N_{\ch\mu}}
    \ar[r, "\text{\'etale}"] &     \sA^{\ch\lambda}
\end{tikzcd}
\]
and the composition of the bottom arrows 
factors through the $(\prod_{\ch\mu} {\mf S}_{N_{\ch\mu}})$-torsor
\[ \textstyle \oo \prod \oo C^{N_{\ch\mu}} \to \oo \prod \oo C^{(N_{\ch\mu})} = \oo C^{\mf R},
\] 
where $\mf S_N$ denotes the symmetric group on $N$ elements.
By induction, we deduce that 
\begin{equation} \label{e:monodromyS}
\source{intersection-complexes-unramified-l-factors}
     \Scr L^{\mf R}[ \abs{\mf R} ]|^*_{\oo\prod \oo C^{N_{\ch\mu}} } 
\cong ( \bt_{\ch\mu} (V_{X,\ch\mu} \ot \IC_C)^{\bt N_{\ch\mu}} )|^*_{\oo \prod \oo C^{N_{\ch\mu}} }. 
\end{equation}
There is a natural $\mf S_{N_{\ch\mu}}$-equivariant structure 
on $(V_{X,\ch\mu} \ot \IC_C)^{\bt N_{\ch\mu}}$ 
compatible with the $\mf S_{N_{\ch\mu}}$-action on $C^{N_{\ch\mu}}$.
On the other hand, we have the map
\[ (\olsf Y^{\ch\mu} \ttimes C)^{\oo\xt N_{\ch\mu}} \to 
\sY^{N_{\ch\mu}\cdot\ch\mu} \xt_{\sA^{N_{\ch\mu}\cdot\ch\mu}} \oo C^{(N_{\ch\mu})} 
\] 
which is a $\mf S_{N_{\ch\mu}}$-torsor, where $\mf S_{N_{\ch\mu}}$
acts on the left hand side in the natural way. 
Now from the definition of $V_{X,\ch\mu}$ we deduce
that the isomorphism \eqref{e:monodromyS} must 
intertwine the $\prod \mf S_{N_{\ch\mu}}$-structures.
By Galois descent we conclude that
$\Scr L^{\mf R} [ \abs{\mf R} ]\cong (\bigotimes_{\ch\mu} 
\Sym^{N_{\ch\mu}}(V_{X,\ch\mu})) \ot \IC_{\oo C^{\mf R}}$. 
Since $\bar\iota^{\mf R} : C^{\mf R} \to \sA^{\ch\lambda}$ is
the normalization of the closure of the stratum $\oo C^{\mf R}$ in $\sA^{\ch\lambda}$,
the middle extension of $\IC_{\oo C^{\mf R}}$ is
isomorphic to $\bar\iota^{\mf R}_!(\IC_{C^{\mf R}})$.

When $X$ is defined over a finite field $\mbb F$ and $k$ is the algebraic closure, the Galois group of $k$ over $\mbb F$ acts naturally on the set of components $\mf B^+_X := \bigcup_{\ch\mu\in \mf c_X} \mB_{X,\ch\mu}$, and the isomorphisms above are clearly equivariant, taking into account that $\IC_{C^{\mf R}}$ is the constant sheaf $ \ol\bbQ_\ell(\frac{|\mf R|}{2})[|\mf R|]$.
\end{proof}




\subsection{Critical dimension}



We can now give a more precise description of the diagonal contribution
$V_{X,\ch\lambda}$ defined in \eqref{e:defVX}, that is, the free vector space on the set $\mB_{X,\ch\lambda}$ of components of critical dimension in the ``open strata'' of the central fiber $\msf Y^{\ch\lambda}$. 

\begin{prop}
\source{intersection-complexes-unramified-l-factors}
\notready \label{prop:Vbasis} 
\source{intersection-complexes-unramified-l-factors}
For $\ch\lambda\in \mathfrak c_X \sm 0$, the set $\mB_{X,\ch\lambda}$ of components of critical dimension on the central fiber $\msf Y^{\ch\lambda}$ consists of
\begin{enumerate}
\item the irreducible components of $\sY^D_{X^\bullet}\cap \msf Y^{\ch\lambda}$ of dimension $\frac 1 2(\len(D)-1)$, 
for $D\in \mathbb N^{\mathcal D}$ with $\varrho_X(D) =\ch\lambda$;
\item the irreducible components of $\msf S^{\ch\lambda} \cap \Gr_G^{\ch\theta}$, 
for $\ch\theta \in \Cal D^G_{\mathrm{sat}}(X)$, embedded in $\msf Y^{\ch\lambda}$ via \eqref{e:actvY}.
\end{enumerate}
\end{prop}

\begin{rem}
\source{intersection-complexes-unramified-l-factors}
\notready \label{rem:othertheta}
\source{intersection-complexes-unramified-l-factors}
We have MV cycles for every $\ch\theta \in \mf c_X^-$, but only those belonging to $\Cal D^G_{\mathrm{sat}}(X)$ 
contribute to $V_{X,\ch\lambda}$ since those correspond to $\sY^{\ch\lambda,\ch\theta}$ which are
connected components of $\sY^{\ch\lambda}$. We will reserve the term \emph{critical dimension} of 
$\msf Y^{\ch\lambda}$ for the maximal dimensions in the two cases above. 
\end{rem}



\begin{proof}
By definition, an element of $\mB_{X,\ch\lambda}$ is a component $\msf b$ of the central fiber $\msf Y^{\ch\lambda}\cap \mathscr Y$, where $\mathscr Y$ is an irreducible component of the smooth stratum $\sY^{\ch\lambda,\ch\theta}$, for some $\ch\theta \in \Cal D^G_{\mathrm{sat}}(X)\cup\{0\}$, such that 
\begin{equation} \label{e:ss-eq}
\source{intersection-complexes-unramified-l-factors}
    \dim (\msf b) = \frac 1 2 ( \dim (\mathscr Y) -1).
\end{equation}

If $\ch\theta=0$, then $\sY^{\ch\lambda,0}$ is the disjoint union of 
connected components $\sY^D_{X^\bullet}$ 
for $D\in \mathbb N^{\mathcal D}$ with $\varrho_X(D) =\ch\lambda$, by Lemma~\ref{lem:oYconnected}. 
Then 
\eqref{e:ss-eq} becomes $\dim(\msf b)
= \frac 1 2 (\len(D)-1) $. 

If $\ch\theta \ne 0$, by Corollary \ref{cor:Ycomp2} the connected components of $\sY^{\ch\lambda,\ch\theta}$ are
in bijection with the closures of 
\[ \sY^{D}_{X^\bullet} \oo\xt \sY^{\ch\theta,\ch\theta} \into \sY^{\ch\lambda,\ch\theta} \]
for $D\in \mbb N^{\Cal D}$ such that $\varrho_X(D) = \ch\lambda-\ch\theta$. The statement now follows from the following lemma, which we write separately, for later use, because it applies to arbitrary $\ch\theta \ne 0$.







\end{proof}


\begin{lem}
\source{intersection-complexes-unramified-l-factors}
\notready \label{lem:centralfibertheta}
\source{intersection-complexes-unramified-l-factors}
For any $\ch\theta \in \mf c_X^- \sm 0$, $D\in \mbb N^{\Cal D}$
and $\ch\lambda = \varrho_X(D) +\ch\theta$, if we denote by $\mathscr Y$ the closure of the image of 
\[\sY^{D}_{X^\bullet} \oo\xt \sY^{\ch\theta,\ch\theta} \into \sY^{\ch\lambda,\ch\theta},
\]
then we have $\dim (\msf Y^{\ch\lambda}\cap \mathscr Y) \le \frac 1 2\len(D) = \frac 1 2(\dim \mathscr Y -1)$. The irreducible components of $\msf Y^{\ch\lambda,\ch\theta}$ for which this is an equality are precisely 
the MV cycles in $\msf S^{\ch\lambda}\cap \ol\Gr^{\ch\theta}_G$, 
embedded in $\msf Y^{\ch\lambda}$ via \eqref{e:actvY}.
\end{lem}
\begin{proof}
Proposition~\ref{prop:centralstrata} implies that $\dim(\msf Y^{\ch\lambda} \cap \mathscr Y)
\le \frac 1 2(\dim \mathscr Y -1)$. 
Since $\ch\theta\ne 0$, we have $\sY^{\ch\theta,\ch\theta}_\red=C$,  and by Lemma \ref{lem:oYconnected} this inequality translates to 
$\dim(\msf Y^{\ch\lambda} \cap \mathscr Y) \le \frac 1 2 \len(D)$. 

If $v\in \abs C$ is the point we are taking central fibers with respect to,
then $\msf Y^{\ch\lambda}$ maps to 
the substack $\sM_{X,v}\subset \sM_{X}$ of maps that are only $G$-degenerate at $v$. 
Recall from Theorem \ref{thm:comp} that $\sM_{X,v}^{\ch\theta}$ 
is contained in the image of 
\[ 
    \act_{\sM,v} : \Bun_{H} \ttimes \Gr_{G}^{\ch\theta} \to \sM_{X,v} 
\]
and the map is birational onto its image. 
We deduce from \eqref{e:Z0act-diagram} and 
Proposition~\ref{prop:Yact-strat} that the fiber product 
$(\Bun_{H} \ttimes \Gr^{\ch\theta}_G) \xt_{\sM_{X}} \msf Y^{\ch\lambda}$ 
has a stratification by 
\[ \bigcup_{\ch\nu} \msf Y^{\ch\lambda-\ch\nu,0} \ttimes (\msf S^{\ch\nu} \cap \Gr^{\ch\theta}_G) 
\] 
where $\ch\nu$ ranges over the weights of $V^{\ch\theta}$.  

Observe that we have an embedding $\ch\Lambda^\pos_G \into \mathbb N^{\Cal D}$
by sending a simple coroot $\ch\alpha \mapsto D_\alpha^+ + D_\alpha^-$. 
By restricting to $X^\circ P_\alpha$ we can deduce that the image of
$\ch\alpha$ in $\pi_1(H)\ot\bbQ$ under \eqref{e:color-inj-piH} is zero. 
Thus, the commutativity of diagram \eqref{e:Y0toBunH} ensures that if
we restrict to the connected component of $\Bun_H$ corresponding to 
$\sY^{D}_{X^\bullet}$, then the stratification above becomes 
\[ \bigcup_{\ch\nu} (\msf Y^{\ch\lambda-\ch\nu}\cap \sY^{D -(\ch\nu-\ch\theta)}_{X^\bullet}) \ttimes (\msf S^{\ch\nu} \cap\Gr^{\ch\theta}_G ),
\]
where $\ch\nu-\ch\theta\in \ch\Lambda^\pos_G \subset {\mbb N}^{\Cal D}$.
In particular, the dimension of the stratum corresponding to $\ch\nu$ is 
\[ \le \frac 1 2 \Bigl(\len(D-(\ch\nu-\ch\theta))-1 \Bigr)+ \brac{\rho_G, \ch\nu-\ch\theta}
\le \frac 1 2 (\len(D)-1) \]
by Proposition~\ref{prop:centralstrata} \emph{unless} $D= \ch\nu-\ch\theta$. 
Thus, in order for $\dim \msf Y^{\ch\lambda,\ch\theta}=\frac 1 2 \len(\ch\lambda-\ch\theta)$, we must have
$\ch\lambda=\ch\nu$ is a weight of $V^{\ch\theta}$ 
and $\msf Y^{\ch\lambda,\ch\theta}$ is birational to an irreducible
component of $\pt \ttimes (\msf S^{\ch\lambda}\cap \Gr^{\ch\theta}_G)$, i.e., 
a Mirkovi\'c--Vilonen cycle. By Lemma~\ref{lem:MVcentralfiber}, 
this latter case always occurs. 
\end{proof}


\begin{rem}
\source{intersection-complexes-unramified-l-factors}
\notready
 By Proposition \ref{prop:Vbasis}, the irreducible components of central Zastava fibers of critical dimension, which give rise to the ``new'' contributions $V_{X,\ch\lambda}$ to the pushforward of the IC sheaf by Proposition \ref{prop:format}, are of two different kinds: those associated to the Zastava space of the open $G$-orbit $X^\bullet$, and those associated to certain strata of the affine Grassmannian. On the other hand, Theorem \ref{thm:heckeIC} gives a similar description of the intersection complex of the global model in terms of the Hecke action on the intersection complex of the global model for $X^\bullet$. 
These two descriptions ``match'' under the nearby cycles functor of Theorem~\ref{thm:nearby}
and the Hecke action on Drinfeld's compactification $\ol\Bun_{N^-}$ (cf.~\cite{BG}). 
\end{rem}
